\chapter*{前言}
\addcontentsline{toc}{chapter}{前言}
\markboth{前言}{前言}

一个结论怎样从定义中得到，一种方法为什么适用于眼前的问题，学过的知识又怎样支持新的认识，是我编写这份讲义时始终关心的事情。数、式与图形构成了熟悉的起点；沿着它们之间的关系继续追问，便会逐渐进入数学的推理与结构。《从数字到数学》由此展开。

这是一份面向高中生的专题讲义，服务于日常学习，也为复习和进一步探索提供材料。可以配合讲授用于相关内容的首次学习，也可以在已有经验之后重新阅读。我相信，经过适当的准备与引导，高中生能够逐步理解这里的概念、论证和抽象观点。

本书依照定义和推理的依赖组织内容，先建立共同的基础，再研究具体的对象。熟悉的问题可以引出概念，直观可以帮助发现关系，定义与证明则使这些认识获得确切的依据。各专题分别形成完整的结构，概念、性质和方法在其中逐步建立；随着学习推进，前面的知识继续参与后面的研究，使读者能够看清每一步的来处，也能知道它在整个问题中的位置。

方法从研究对象本身的关系与结构中产生。向量、函数以及齐次线性方程的解所具有的线性结构，使线性代数的方法适用于相应的问题；圆与椭圆之间存在仿射关系，一些几何性质才能随之传递。具体解法由此有了可以追溯的依据。掌握这些依据，读者便有机会在题目改变时重新选择表示、组织推导，并得到自己需要的结论。

严谨决定推理是否成立，取舍决定哪些内容值得展开。关键的条件与推理需要交代，较为抽象的内容需要适当准备，直接的推论和常规计算则可以留给读者完成。我把详略的依据放在内容的作用和理解的需要上：一个问题应当讲到依据清楚、关系明确，也应当有恰当的结束位置。作者掌握的知识与工具，用于帮助读者建立认识，并为读者留下继续思考的空间。

本书不应试，但不反应试。课程与考试中的问题值得认真处理，基本运算和常用方法也需要练习。本书把这些问题放在相应的知识体系中，说明方法的来源和条件的作用。对数学思想的理解，应当能够支持具体的解题；解题中的经验，也应当能够整理为更清楚的认识。计算、推理与应用在同一学习过程中相互支持。

\newpage
本书的六章各自围绕一类问题展开。第一章从数列进入极限、级数及迭代，建立研究无限过程的基础；第二章从函数概念出发整理性质，研究连续性；第三章从几何向量进入线性空间、变换和复数，为后面的结构观点准备工具。

第四章以二次曲线为对象，把方程、几何关系和多种表示方法组织在同一体系中；第五章分别从变化率与累积量建立微分、积分，证明二者的联系，再研究应用，最终以微分方程收束；第六章从计数原理建立概率模型，由随机变量的分布走向抽样、估计、检验和回归。前面建立的极限、线性结构和微积分，在新的问题中继续发挥作用。

各章的内容有所承接，后面的讨论将使用前文已经建立的概念与方法。第一章需集合与常用逻辑用语的基础，第二章使用函数的基本经验并承接极限概念，第三章由几何向量展开。第四章主线使用高中代数与几何，星号部分调用第三章的线性代数知识。第五章承接前两章，并在选讲中使用线性思想。第六章以高中代数及函数语言为基础，大数定律调用极限思想，选讲进一步使用前面建立的工具。

正文中的引入用于提出问题，定义确定研究对象，命题与证明建立性质。阅读推导时，可以留意条件在哪一步起作用、所选的表示为何有用，并尝试自行补齐省略的计算。知识与方法清单用于整理和查阅，可以在学习相应内容后反复使用。

星号标记选讲范围和认识层次，题目难度另由问题决定。主线可独立完成相应专题的基本学习，读者可以跳过全部星号内容。星号线通常围绕同一对象深化解释，在各节之间形成连续的联系，又接回主线中的具体问题。读者可以根据基础和需要穿插阅读，也可以完成主线后再集中学习。已有经验往往会在重读时帮助理解新的观点。

例题保留题面，供课堂展开与独立思考。节内习题落实相应内容，章末综合习题作为覆盖全章的总习题集，单题可以考察一个知识点，也可以联合多个知识点。配套习题解另册收录例题与习题的解析，供备课、核对与复习使用。读者可以先独立尝试，再结合解析检查推理；遇到困难时，也可以回到正文寻找定义、性质或相近的问题。阅读解答时，除了核对结果，还可以回顾对象的选择、表示方式和关键依据，比较同一方法在不同问题中的用途。

附录汇集常用记号、不等式及公式，供学习时查阅。

\begin{flushright}
Andrew Elizabeth / EasterGhost\\
2026年11月14日
\end{flushright}
